\chapter{PD-L1 (Immunotherapy) Tests}
\section*{What Is This Test?} PD-L1 testing measures programmed death-ligand 1 expression in tumor tissue to help determine eligibility for selected immunotherapies.\section*{Alternative Names} PD-L1 IHC; programmed death-ligand 1 test; immune checkpoint biomarker.\section*{Principle} Immunohistochemistry detects PD-L1 staining and reports a scoring system such as tumor proportion score or combined positive score.\section*{Why This Test Is Ordered} It supports treatment decisions in selected lung, breast, gastric, esophageal, cervical, urothelial, and other cancers.\section*{Primary vs Secondary Test} It is a secondary pathology biomarker after tissue diagnosis; the assay and drug indication must match.\section*{How This Test Is Performed} A pathologist stains biopsy or surgical tissue using a validated companion or complementary diagnostic assay.\section*{What Preparation Is Needed} No patient preparation beyond tissue sampling.\section*{Understanding Results} Positive and negative thresholds vary by cancer, specimen, assay, and treatment; a negative result does not exclude all benefit.\section*{Follow-up Tests} HER2, hormone receptors, genomic tests, staging imaging, and oncology consultation may follow.\section*{References} \begin{enumerate}\item MedlinePlus. PDL1 (Immunotherapy) Tests.\item American Society of Clinical Oncology and College of American Pathologists. PD-L1 testing guidance.\end{enumerate}\clearpage\section*{Understanding Results and Follow-up}
The laboratory report should identify the specimen, method, units, reference interval or decision limit, and whether the result is positive, negative, abnormal, or inconclusive. Reference intervals vary with age, sex, pregnancy, specimen type, assay, and laboratory; a result outside a range is not automatically a diagnosis, and a result within a range does not always exclude disease. Discuss unexpected results with the ordering clinician, who can decide whether repeat or confirmatory testing, treatment, or referral is appropriate. Do not change medicines or delay urgent care based on an isolated result.
\section*{Regulatory Status and Authorization Date}This chapter describes a test category, not a specific commercial product. FDA, CDSCO, EU IVDR, and MHRA approval or registration dates therefore cannot be assigned without the assay manufacturer, product name, instrument, and intended use. For laboratory-developed tests, an individual agency approval date may not exist. See the Preface for the interpretation of regulatory status.
\clearpage
